\documentclass[]{article}

\usepackage{amsmath}
\usepackage{multirow}
\usepackage{natbib}
\usepackage{graphicx} 


\begin{document}

This paper sought to address a central challenge in empirical growth economics: distinguishing between the factors that endogenously determine the level of investment and those that moderate the speed of convergence to a steady-state. We developed and applied a two-stage empirical framework to a synthetic panel dataset, which was generated from a structural growth model with known parameters, allowing for a controlled validation of our methodology.

Our analysis first confirmed the presence of a significant endogenous investment mechanism. Using a fixed-effects panel regression, we found a positive and statistically significant elasticity of the investment rate with respect to GDP per capita. This result establishes that as economies develop, their propensity to invest increases, creating a crucial feedback loop that drives capital formation. This finding underscores the importance of treating the investment rate not as an exogenous parameter but as a dynamic variable intrinsically linked to the development process itself.

In the second stage, we modeled the dynamics of capital accumulation as a function of the distance to an income-conditioned steady-state. Our results revealed a strong and rapid baseline convergence process, with countries closing approximately 35\% of the gap to their steady-state capital stock each period. The key innovation was the use of an instrumented moderated transition model to test whether policy variables influence this convergence velocity. We found that trade openness acts as a marginal accelerator of capital deepening; more open economies converge to their steady-state faster than their more closed counterparts. In contrast, the share of government expenditure in GDP did not have a statistically significant effect on the speed of convergence within our framework.

From these results, we have learned that it is both possible and insightful to empirically separate the determinants of an economy's long-run equilibrium from the factors that govern the speed of its transition. The framework presented here successfully disentangled the level effect of income on investment from the speed effect of policy on capital accumulation. The findings suggest that policies like trade openness may enhance growth not just by potentially raising the steady-state level of capital, but also by increasing the efficiency of the convergence process itself. This distinction provides a more nuanced understanding of the channels through which income and policy shape economic growth dynamics.

\end{document}
                